\chapter{Metodologija vrednotenja}
\label{ch:metodologija}

Vrednotenje je bilo zasnovano tako, da ločeno obravnava kakovost naučenega
modela in delovanje celotne fizične naprave. Razvojna pravilnost modela sama po
sebi namreč ne pove, ali bodo zanesljivo delovali tudi zvočnik, vgrajeni
mikrofon, predobdelava, nevronski pospeševalnik, odločitveni filter, serijski
diagnostični izhod in zaslon. Po drugi strani uspešen prikaz nekaj izbranih
zvokov še ni neodvisen dokaz splošne pravilnosti modela.

Med razvojem je bilo zaporedoma ovrednotenih več pomembnih izvedb. Rezervirani
preskus je uporabil takrat zamrznjeno različico YAMNet-256. Njegovi
rezultati so razkrili predvsem težave pri kratkotrajnih dogodkih in so bili nato
uporabljeni pri načrtovanju izboljšav. Končna naprava uporablja večjo različico
YAMNet-1024, izboljšan učni nabor, razred drugih zvokov in pragove po razredih.
Najprej smo jo preverili z nadzorovanim sprejemnim preskusom. Ker so bili
njegovi posnetki predhodno uporabljeni med razvojem, ga ne predstavljamo kot
neodvisen preskus posploševanja. Po zadnjih spremembah smo pripravili še nov
končni fizični preizkus s 60 slušno pregledanimi posnetki, katerih viri niso
sodelovali pri učenju. Njegov protokol, vrstni red, programska slika in model
so bili zamrznjeni pred prvim predvajanjem.

\section{Cilji vrednotenja}

Z vrednotenjem smo želeli odgovoriti na štiri praktična vprašanja:

\begin{enumerate}
  \item Ali celotna naprava v nadzorovanih pogojih potrdi pričakovani razred
        zvoka?
  \item Ali govor in raznoliki zvoki razreda drugo povzročijo lažen
        nevarnostni alarm?
  \item Kako hitro po začetku predvajanja sistem prvič potrdi pričakovani
        razred?
  \item Kolikšna sta čas izvajanja obdelovalnih stopenj in pomnilniški odtis
        izvedbe?
\end{enumerate}

Prva tri vprašanja obravnavajo obnašanje sistema od zvočnika do končne
odločitve. Četrto vprašanje obravnava primernost izvedbe za mikrokrmilnik. Pri
tem ločimo čas računanja od zakasnitve zaznave. Vsak vhod obsega približno
\SI{975}{\milli\second} zvoka, nova obdelava pa se zaradi prekrivanja začne
vsakih \SI{480}{\milli\second}. Odločitveni filter lahko za
potrditev zahteva več zaporednih blokov. Čas izvajanja modela zato ni enak času
od začetka zvoka do prikaza opozorila.

\section{Razvojni in ocenjevalni preskusi}

Med izdelavo sistema je bilo izvedenih več vrst preskusov. Njihov namen in
dokazna vrednost nista enaka. Razvojni podatki so bili uporabljeni za izbiro
modela, nastavitev pragov in odpravljanje napak. Rezervirani podatki so bili
odprti šele po zamrznitvi starejše izvedbe. Sprejemni preskus končne naprave pa
je uporabil že pregledane kanonične primere, da je bilo mogoče preveriti
delovanje točno tiste različice, ki je ostala na ploščici.

\begin{table}[!ht]
  \centering
  \caption{Vloga posameznih stopenj vrednotenja.}
  \label{tab:stopnje-vrednotenja}
  \small
  \begin{tabular}{p{0.23\textwidth}p{0.60\textwidth}}
    \toprule
    Stopnja & Namen in dovoljena razlaga \\
    \midrule
    Razvojno preverjanje modela & Izbira učnega postopka in primerjava različic
    YAMNet-1024; končna sedemrazredna različica uporablja 413 razvojnih
    posnetkov in ni končna ocena. \\
    Rezervirani fizični preskus & Preverjanje takrat zamrznjene izvedbe
    YAMNet-256 na 70 prej neuporabljenih virih; rezultat velja samo za to
    izvedbo. \\
    Kalibracija in dimni preskusi & Nastavitev vhodnega praga in pragov po
    razredih ter preverjanje zagona, zaslona in vseh sistemskih izhodov;
    rezultati niso uporabljeni kot ocena pravilnosti. \\
    Nadzorovani sprejemni preskus & Preverjanje končne naprave YAMNet-1024 na
    odobrenih primerih in pri treh ravneh predvajanja; ne meri posploševanja na
    nove posnetke. \\
    Neodvisni končni fizični preizkus & Primarna ocena zamrznjene končne
    naprave na 60 novih virih, ločenih od učenja; vključuje pravilne zaznave,
    prevladujoči izhod, lažne nevarnostne alarme in odzivni čas. \\
    \bottomrule
  \end{tabular}
\end{table}

Razvojno preverjanje končnega modela je obsegalo 413 posnetkov, ki niso
sodelovali pri učenju zaključne plasti. V njem je bilo tudi 135 primerov
razreda drugo. Rezultat je bil uporabljen pri izbiri med razvojnimi
različicami, zato ga obravnavamo kot razvojni rezultat.
Pred fizičnimi preskusi so bili opravljeni tudi slušni pregledi izrezov in
povečanih različic. Pri strelu in razbitju stekla smo preverili, ali je dogodek
res prisoten in ali ni odrezan na začetku ali koncu. Pri nevihti smo zavrnili
posnetke z izrazitim visokim tonom v ozadju. S tem smo preverjali kakovost
dražljajev, ne odziva naprave.

Dimni preskus pomeni kratek funkcijski pregled po spremembi programske opreme
ali modela. Z njim smo preverili, ali se naprava zažene, ali delujeta mikrofon in
zaslon ter ali je mogoče sprožiti vse pričakovane izhode. Posnetki in rezultati
takšnih preskusov so lahko vplivali na naslednjo razvojno odločitev, zato niso
vključeni v končne deleže uspešnosti. Enako velja za nedokončano primerjalno
predvajanje starega nabora 70 zvokov na novi izvedbi.

\section{Rezervirani preskus predhodne izvedbe}
\label{sec:metodologija-rezervirani}

Rezervirani preskus z oznako
\texttt{STM32N6-HAZARD6-FINAL-001} in poskusom \texttt{FE-A01} je bil
pripravljen za takrat končno različico YAMNet-256. Pred prvim zajemom so bili
zamrznjeni model, programska oprema, pragovi, predobdelava, seznam dražljajev,
njihov vrstni red in pravila ocenjevanja. Med 70 poskusi ni bila ponovljena
nobena neuspešna napoved in nastavitve niso bile spremenjene.

\subsection{Sestava preskusnega nabora}

Preskus je vseboval 60 pozitivnih in 10 zahtevnih negativnih primerov. Za vsak
izmed šestih izhodov je bilo izbranih deset različnih izvornih posnetkov. Med
izhodi so bili pasji lajež, razbitje stekla, strel oziroma streljanje, sirena,
govor in nevihta. Pozitivni posnetki so izvirali iz ocenjevalnega dela zbirke
Freesound Dataset 50K, ki med razvojem te izvedbe ni bil uporabljen. Negativni
posnetki so izvirali iz petega dela zbirke Environmental Sound Classification
50. Vsebovali so ploskanje, trkanje po lesenih vratih, jok dojenčka, tiktakanje
ure, dež, motorno žago, prasketanje ognja, avtomobilsko hupo, kašljanje in
ognjemet. To niso bili tihi kontrolni primeri, temveč zvoki, ki bi se lahko
zamenjali z enim od ciljnih razredov.

Izbira je bila določena s pravili nad metapodatki in ni upoštevala napovedi
modela ali kasnejšega poslušanja rezultatov. Pri razbitju stekla je bila
zahtevana oznaka, ki pomeni dejanski lom, pri nevihti pa oznaka za grmenje.
Metapodatkovno označeni umetni streli so bili izločeni. Govor je obsegal štiri
moške, štiri ženske in dva otroška posnetka. Izločeni so bili igrani, kričani,
šepetani, peti, umetno ustvarjeni ali višinsko spremenjeni glasovi. Posnetki
istega ciljnega razreda so izvirali od različnih nalagalcev, s čimer smo
zmanjšali odvisnost od ene snemalne seje.

\subsection{Priprava dražljajev}

Vsi viri so bili pretvorjeni v enokanalni zapis s frekvenco vzorčenja
\SI{44.1}{\kilo\hertz} in šestnajstbitnimi celoštevilskimi vzorci. Po
normalizaciji je bil samodejno izbran najmočnejši ustrezni izsek. Neprekinjeni
zvoki so uporabljali petsekundske enote, kratkotrajni dogodki pa dve sekundi
dolge enote z dogodkom v sredini. Enote so bile z vnaprej določenimi premori
ponovljene do skupne dolžine najmanj deset sekund.

Osnovni cilj normalizacije je bila največja efektivna vrednost v oknu dolžine
\SI{100}{\milli\second}, in sicer \SI{-12}{\decibel} glede na polno digitalno območje, ter najvišji vrh
\SI{-1}{\decibel} glede na polno digitalno območje. Zaradi kratkosti dogodkov
je bilo razbitje stekla nato ojačeno za \SI{3}{\decibel}, strel pa za
\SI{10}{\decibel}. Če bi ojačenje povzročilo omejevanje več kot petih odstotkov
vzorcev, ga je pripravljalni program zmanjšal samo toliko, da je ostal pod to
mejo. Za negativne primere so bile že pred preskusom določene udobne absolutne
ravni predvajanja, saj se njihove izvorne glasnosti močno razlikujejo.

Za vsak dražljaj so bile shranjene izvorna in izhodna kriptografska kontrolna
vsota, uporabljeno ojačenje, položaj izseka, število ponovitev in delež
omejevanja. Tako je bilo mogoče preveriti, da se med pripravo in predvajanjem
datoteka ni spremenila.

\subsection{Fizični pogoji in izvedba}

Glasnost sistema Windows je bila nastavljena na \SI{75}{\percent}, razdalja med
zvočnikom in vgrajenim mikrofonom pa na \SI{30}{\centi\meter}. Razvojna plošča
in zvočnik sta med poskusom ostala na istem mestu v mirnem prostoru. Pred
začetkom je bil enkrat pritisnjen gumb za ponastavitev, med preskusom pa
ponastavljanje ni bilo dovoljeno.

Slika \ref{fig:preskusna-postavitev} prikazuje uporabljeno namizno postavitev.
Merilni trak je bil uporabljen za ponovljivo nastavitev razdalje med zvočnikom
in mikrofonom na razvojni plošči. Fotografija dokumentira razporeditev opreme,
ne pa posameznega rezultata preizkusa.

\begin{figure}[H]
  \centering
  \includegraphics[width=0.68\textwidth,trim=150 0 100 700,clip]{20260817_preskusna_postavitev.jpg}
  \caption{Preskusna postavitev z zvočnikom, merilnim trakom in razvojno
  ploščo STM32N6570-DK.}
  \label{fig:preskusna-postavitev}
\end{figure}

Vsak zajem se je začel tri sekunde pred predvajanjem in je trajal 20 sekund.
Diagnostični zapisi so bili zajeti prek serijske povezave na vratih COM3 s
hitrostjo 14.400 bitov na sekundo. Vrstni red je bil sestavljen tako, da dva
zaporedna poskusa nista imela istega pričakovanega izhoda. Program je izvajalcu
prikazoval le oznako poskusa in napredek, napovedi pa je shranjeval samodejno.
Tehnična napaka bi dovolila nadomestni zajem samo ob ohranitvi prvotnega
neuporabnega zapisa; napačna napoved sama po sebi ni bila razlog za ponovitev.

\section{Sprejemni preskus končne naprave}
\label{sec:metodologija-sprejemni}

Končna izvedba YAMNet-1024 je bila ovrednotena z nadzorovanim sprejemnim
preskusom \texttt{STM32N6-HAZARD6-ACCEPTANCE-001}. Izveden je bil enkrat kot
poskus \texttt{AT-A01}. Pred začetkom so bili zamrznjeni model, podpisana
programska slika, odločitveni pragovi, manifest in vrstni red predvajanja.
Identiteta modela, programske slike in vsake zvočne datoteke je bila preverjena
z 256-bitno kriptografsko kontrolno vsoto.

Preskus je uporabil devet predhodno poslušanih in odobrenih izvornih posnetkov.
Za vsakega izmed šestih modelskih izhodov je bil izbran en jasen kanoničen
primer. Vsak je bil predvajan pri nominalni ravni ter po utišanju za
\SI{6}{\decibel} in \SI{12}{\decibel}. S tem je nastalo 18 pozitivnih poskusov.
Dodani so bili še trije varni negativni poskusi: dež, veter in trkanje po
lesenih vratih. Ti so bili predvajani samo pri nominalni ravni, saj bi jih z
dodatnim utišanjem spremenili v prelahke negativne primere. Skupaj je preskus
vseboval 21 poskusov, kot prikazuje tabela \ref{tab:sprejemni-nabor}.

\begin{table}[!ht]
  \centering
  \caption{Sestava nadzorovanega sprejemnega preskusa končne naprave.}
  \label{tab:sprejemni-nabor}
  \begin{tabular}{lrrr}
    \toprule
    Skupina & Nominalna raven & \SI{-6}{\decibel} & \SI{-12}{\decibel} \\
    \midrule
    Šest modelskih izhodov & 6 & 6 & 6 \\
    Trije varni zvoki & 3 & 0 & 0 \\
    \midrule
    Skupaj & 9 & 6 & 6 \\
    \bottomrule
  \end{tabular}
\end{table}

Semantični pregled je bil pri kratkotrajnih dogodkih posebej pomemben.
Posnetek razbitja stekla je moral vsebovati jasno in celovito razbitje, ne le
rokovanja s steklom, trka ali cingljanja. Posnetek strela je vseboval jasen
posamezen strel, ki je bil v dražljaju štirikrat ponovljen. Posnetek nevihte ni
vseboval prej opaženega motečega visokega tona. Ta nadzor zmanjšuje nevarnost,
da bi napravo ocenjevali z napačno označenim ali vsebinsko neustreznim zvokom.

Glasnost sistema Windows je bila nastavljena na \SI{70}{\percent}, razdalja med
zvočnikom in mikrofonom pa na \SI{30}{\centi\meter}. Položaj naprav se ni
spreminjal, prostor je bil čim bolj tih, ploščica pa je bila pred začetkom
enkrat ponastavljena. Vrstni red je bil določen s semenom 811. Pred vsakim
predvajanjem je program tri sekunde zajemal stanje prostora, nato pa je
samodejno predvajal datoteko in shranil serijske zapise. Neuspešna napoved ni
bila ponovljena.

Pri pozitivnem poskusu je bil rezultat uspešen, če se je pričakovani razred med
predvajanjem vsaj enkrat pojavil kot potrjena odločitev. Pri varnem zvoku je bil
uspeh odsotnost potrjenega nevarnostnega razreda; izhodi \emph{čakanje},
\emph{neznano} in \emph{govor} so bili dovoljeni. Primarni rezultat je delež
uspešnih devetih poskusov pri nominalni ravni. Poskusi pri znižanih ravneh so
sekundarno preverjali občutljivost na glasnost. Ker je bila vsaka kombinacija
izvornega posnetka in ravni predvajana samo enkrat, rezultate prikazujemo
opisno in jih ne razlagamo kot oceno pravilnosti na celotni populaciji zvokov.

\section{Neodvisni končni fizični preizkus}
\label{sec:metodologija-koncni}

Po dokončanju razreda drugo, pragov po razredih in polovičnega prekrivanja
vhodnih oken smo pripravili novi preizkus z oznako
\path{STM32N6-HAZARD6-FINAL-003}. Edini izvedeni poskus ima oznako
\path{FE4-A01}. Pred začetkom so bili zamrznjeni model, uteži, podpisana
aplikacija, seznam dražljajev, njihov vrstni red in pravila analize. Manifest
ima 256-bitno kriptografsko kontrolno vsoto
SHA-256:
\begin{center}
  \small\texttt{0355fca0f9a3d26be1e41b4fd9896f7}\newline
  \texttt{fb7c9f567557a7dde23826dbb36c1342b}.
\end{center}

\subsection{Izbira in pregled posnetkov}

Pred zamrznitvijo je bil opravljen slepi slušni pregled 118 kandidatov. Med
poslušanjem napovedi modela in ploščice niso bile prikazane. Kandidat je bil
primeren samo, če je vseboval običajen in jasno prepoznaven dogodek, če ta ni
bil odrezan na začetku ali koncu, če je bila slišnost normalna ter če ni bil
označen zvočni artefakt. Identitete virov so bile preverjene tudi med
različnimi zbirkami. En kandidat pasjega laježa z istim izvornim posnetkom kot
v učni množici je bil zato izločen.

Iz primernih kandidatov je bilo z vnaprej določenima semenoma izbranih 60
različnih virov, po deset za vsak sistemski razred. Razredi so bili pasji
lajež, razbitje stekla, strel, drugo, sirena in govor. Nevihta ni samostojen
sistemski razred, saj se njen modelski izhod pred odločanjem prišteje razredu
drugo. Deset primerov razreda drugo predstavlja deset različnih vrst zvoka:
alarm, aplavz, zvonec, ploskanje, računalniško tipkovnico, pok, posodo in
kuhinjske pripomočke, ognjemet, trkanje ter veter. Izbira ni uporabljala
napovedi končnega modela.

Vrstni red je bil naključno določen pod pogojem, da zaporedna poskusa nimata
istega pričakovanega razreda. Vsak dražljaj je bil normaliziran na največjo
efektivno vrednost v \SI{100}{\milli\second} oknu približno
\SI{-42}{\decibel} glede na polno digitalno območje. Sistem Windows je ostal
na \SI{100}{\percent} glasnosti, zvočnik pa je bil \SI{30}{\centi\meter} od
vgrajenega mikrofona. Položaj naprav in prostor se med preizkusom nista
spreminjala.

\subsection{Izvedba in merila}

Vsak zajem je trajal 15 sekund. Prve tri sekunde so bile namenjene zajemu
stanja pred predvajanjem, nato je program samodejno predvajal zamrznjeno
datoteko in shranil serijski dnevnik. Za primarno merilo smo uporabili samo
časovno območje od začetka predvajanja naprej. Poskus je bil uspešen, če je
sistem v tem območju vsaj enkrat potrdil pričakovani razred. Napačna napoved ni
bila razlog za ponovitev.

Poleg primarnega merila smo za vsak poskus določili prevladujoči potrjeni izhod
in čas od začetka predvajanja do prve pravilne potrditve. Prevladujoči izhod je
razred z največ potrjenimi okvirji; ob izenačenju odloča večja najvišja ocena.
Lažni nevarnostni alarm smo merili pri govoru in razredu drugo. Pomeni vsaj en
potrjeni izhod pasjega laježa, razbitja stekla, strela ali sirene po začetku
predvajanja. Posebej smo prešteli poskuse z več kot enim takim okvirjem ter
izhode, ki so nastali v treh sekundah pred predvajanjem. Tako preneseno stanje
prejšnjega poskusa ni bilo napačno pripisano novemu dražljaju.

Za skupni in razredne deleže smo izračunali Wilsonovo
petindevetdesetodstotno območje zaupanja. Ker je v vsakem razredu le deset
poskusov, so območja široka. Rezultat je zato ocena celotnega sistema v
določeni postavitvi zvočnik--mikrofon, ne splošna pravilnost za vse resnične
zvoke.

\section{Zajem podatkov in merila uspešnosti}

Vdelana programska oprema je za vsak obdelani blok v serijski diagnostični
izhod zapisala stanje vhodne aktivnosti, šest izhodnih ocen, tri najvišje
razrede, potrjeno odločitev in čase posameznih stopenj. Zajemni program je
vsakemu poskusu dodal oznako modela, programske različice, pričakovani razred,
nastavljeno glasnost, razdaljo, kontrolno vsoto dražljaja in pot do izvornega
zapisa. Izvirni zapisi so bili ohranjeni, analizni program pa je iz njih izdelal
preglednice in grafe.

Pri analizi smo uporabljali naslednja merila:

\begin{itemize}
  \item \emph{potrjena pravilna zaznava} pomeni, da je odločitveni filter vsaj
        enkrat potrdil pričakovani razred;
  \item \emph{pričakovani razred na prvem mestu} pomeni, da je imel ta razred
        v vsaj enem aktivnem bloku najvišjo glajeno oceno, tudi če ni dosegel
        pogojev za potrditev;
  \item \emph{aktivni blok} je blok, pri katerem je energijsko merilo preseglo
        vhodni prag aktivnosti;
  \item \emph{lažen nevarnostni alarm} pomeni, da je sistem pri govoru ali
        razredu drugo potrdil enega od štirih nevarnostnih razredov;
  \item \emph{delež uspešnosti} je razmerje med številom uspešnih in številom
        vseh poskusov v obravnavani skupini.
\end{itemize}

Notranja ocena modela je shranjena tudi kot najvišja dosežena ocena
pričakovanega razreda. Te vrednosti ne imenujemo verjetnost pravilnosti, saj
model ni bil posebej umerjen tako, da bi na primer ocena 0,8 v vseh pogojih
pomenila osemdesetodstotno verjetnost pravilne odločitve.

Pri rezerviranem in neodvisnem končnem preskusu je bilo za deleže izračunano
petindevetdesetodstotno območje zaupanja po Wilsonu
\cite{wilson1927interval}. Za opaženi delež
$\hat{p}=k/n$, število poskusov $n$ in vrednost $z=1{,}96$ sta meji

\begin{equation}
  \frac{\hat{p}+\frac{z^2}{2n}
  \mathbin{\pm}z\sqrt{\frac{\hat{p}(1-\hat{p})}{n}+\frac{z^2}{4n^2}}}
  {1+\frac{z^2}{n}}.
\end{equation}

Pri samo desetih poskusih na razred je takšno območje široko. Zato razlike med
razredi opisujemo kot opažanje v danem preskusu in ne kot natančno oceno
splošne pravilnosti. Za sprejemni preskus območja zaupanja po razredih niso
smiselna, saj tri ravni izvirajo iz istega posnetka in zato niso tri neodvisno
izbrani primeri.

\section{Merjenje časa izvajanja in pomnilnika}

Čas predobdelave, sklepanja in naknadne obdelave meri vdelana programska oprema
s števcem ciklov procesorskega jedra. Število ciklov se glede na frekvenco jedra
pretvori v milisekunde in zapiše ob posameznem obdelanem bloku. Analiza izloči
zapise, pri katerih katera izmed časovnih vrednosti manjka, ter poroča število
popolnih zapisov in povprečje. Gre za telemetrijo same naprave, ne za neodvisno
meritev z osciloskopom ali zunanjim signalom. Prav tako ta čas ne vključuje
čakanja, da se napolni skoraj enosekundno vhodno okno, niti morebitnega čakanja
na potrditev v več zaporednih blokih.

Pomnilniški odtis je bil določen iz povezovalnega zemljevida programa, poročila
prevajalnika nevronskega modela in velikosti ustvarjenih binarnih datotek.
Ločeno obravnavamo notranji pomnilnik aplikacije, notranje območje začasnih
aktivacij nevronskega pospeševalnika, zunanji pomnilnik HyperRAM za dva
zaslonska medpomnilnika ter particiji zunanjega bliskovnega pomnilnika za
aplikacijo in uteži. Te velikosti pripadajo različnim fizičnim območjem, zato
njihovih odstotkov ne seštevamo. Velikost datoteke z razhroščevalnimi podatki
prav tako ni enaka velikosti slike, ki se zapiše na napravo.

Električna moč in poraba energije nista bili izmerjeni. V nalogi zato ne
podajamo številčne trditve o energijski učinkovitosti, čeprav je izvajanje na
namenskem nevronskem pospeševalniku pomemben del zasnove.

\section{Veljavnost in omejitve}

Rezervirani preskus predhodne izvedbe ima dobro ločitev med razvojnimi in
ocenjevalnimi viri, vendar njegovi rezultati ne opisujejo pozneje izdelanega
modela YAMNet-1024. Sprejemni preskus uporablja le devet že poznanih izvornih
posnetkov. Neodvisni končni fizični preizkus uporablja nove vire, ločene od
učenja, in zato predstavlja glavni rezultat končne naprave. Tudi ta preskus pa
zaradi omejenega števila virov ne dokazuje splošne pravilnosti v poljubnem
okolju.

Fizični preskusi so bili izvedeni z eno razvojno ploščico, njenim
vgrajenim mikrofonom, enim zvočnikom, eno razdaljo in v enem prostoru.
Odstotek glasnosti sistema Windows ter digitalna raven v decibelih ne
predstavljata umerjene ravni zvočnega tlaka pri mikrofonu. Na rezultat lahko
vplivajo frekvenčni odziv zvočnika, odboji v prostoru, zunanji šum in trenutna
poravnava kratkega dogodka z obdelovalnim oknom. Še posebej pri strelu in
razbitju stekla lahko majhen časovni premik spremeni vsebino posameznega
bloka.

Nabor razreda drugo z desetimi kategorijami je širši od treh varnih primerov
sprejemnega preskusa, vendar še vedno ne zajema vseh možnih motečih zvokov. Za
močnejšo oceno v resničnem okolju bi potrebovali več novih virov in ponovitev,
umerjene ravni zvoka, različne prostore in razdalje ter večurno opazovanje
lažnih alarmov.
